\title{Chain-of-Thought Prompting Elicits Reasoning in Large Language Models}

\begin{abstract}
We investigate how producing a \textit{chain of thought}—a sequence of intermediate reasoning steps—can greatly enhance the capability of large language models to perform complex reasoning tasks. Specifically, we demonstrate that such reasoning skills naturally arise in sufficiently large language models through a straightforward approach called \textit{chain-of-thought prompting}, where a few examples of chains of thought are included as exemplars in the prompt.

Evaluations on three large language models reveal that chain-of-thought prompting boosts performance across a variety of arithmetic, commonsense, and symbolic reasoning challenges. The improvements observed are substantial. For example, using only eight chain-of-thought exemplars to prompt a PaLM 540B model achieves state-of-the-art results on the GSM8K benchmark for math word problems, outperforming even finetuned GPT-3 equipped with a verifier.
\end{abstract}

\section{Introduction}
Recent advancements in NLP have been transformed by language models \citep[][\textit{inter alia}]{peters-etal-2018-deep,devlin-etal-2019-bert,brown2020language}. Increasing the size of language models has been shown to yield numerous advantages, including better performance and enhanced sample efficiency \citep[\textit{inter alia}]{kaplan2020scaling,brown2020language}. Nonetheless, merely enlarging model size has not sufficed to reach high accuracy on challenging problems such as arithmetic, commonsense, and symbolic reasoning \citep{rae2021scaling}.

In this paper, we examine how the reasoning capabilities of large language models can be unlocked through a simple strategy inspired by two insights. First, methods for arithmetic reasoning improve when models generate natural language explanations leading to the final solution. Previous work has enabled models to produce natural language intermediate steps by either training from the ground up \citep{ling-etal-2017-program} or finetuning pretrained models \citep{cobbe2021training}, alongside neuro-symbolic approaches that utilize formal languages instead of natural language \citep{roy-roth-2015-solving, chiang-chen-2019-semantically,amini-etal-2019-mathqa, chen2019neural}. Second, large language models hold great promise for in-context few-shot learning through \emph{prompting}. Instead of finetuning a separate model for each new task, one can simply provide a few input-output exemplars within the prompt that demonstrate the task. Notably, this approach has been successful for various straightforward question-answering tasks \citep{brown2020language}.

Nevertheless, both of the aforementioned approaches have significant drawbacks. For rationale-augmented training and fine-tuning techniques, generating a substantial collection of high-quality rationales is expensive and considerably more complex than the straightforward input--output pairs typically employed in conventional machine learning. Meanwhile, the traditional few-shot prompting strategy utilized in \citet{brown2020language} tends to perform poorly on tasks demanding reasoning capabilities and often shows limited improvement even as the language model's size increases \citep{rae2021scaling}. In this study, we integrate the advantages of these two concepts while circumventing their respective limitations. Specifically, we investigate the capacity of language models to execute few-shot prompting for reasoning challenges by using prompts composed of triples: $\langle$input, \emph{chain of thought}, output$\rangle$.

A \emph{chain of thought} denotes a sequence of intermediate natural language reasoning steps that culminate in the final answer; we designate this method as \textit{chain-of-thought prompting}. An example of such a prompt is provided in \cref{fig:pull-figure}.

We conduct experimental evaluations across arithmetic, commonsense, and symbolic reasoning benchmarks, demonstrating that chain-of-thought prompting consistently surpasses conventional prompting, occasionally by a substantial margin. \cref{fig:pull-bar-chart} depicts one such outcome—on the GSM8K dataset of math word problems \citep{cobbe2021training}, chain-of-thought prompting with \palm{} 540B considerably outperforms standard prompting and sets new state-of-the-art results. The significance of a prompting-only approach lies in its independence from large training datasets and its ability to enable a single model checkpoint to tackle multiple tasks without sacrificing generality. This research highlights how large language models can effectively learn from a few examples using natural language task descriptions (in contrast to relying on extensive training datasets to implicitly learn input-output patterns).

\section{Chain-of-Thought Prompting} When tackling a complex reasoning task, such as a multi-step math word problem, people often reflect on their own thought process. Typically, they break the problem down into intermediate steps and address each step before arriving at the final solution: \textit{``After Jane gives 2 flowers to her mom she has 10 $\ldots$ then after she gives 3 to her dad she will have 7 $\ldots$ so the answer is 7.''} The objective of this work is to enable language models to produce a comparable \textit{chain of thought}—a logical sequence of intermediate reasoning steps culminating in the solution. We demonstrate that sufficiently large language models can generate such chains of thought when examples illustrating chain-of-thought reasoning are included in the few-shot prompting exemplars.

\cref{fig:pull-figure} illustrates a model generating a chain of thought to solve a math word problem that it would have otherwise answered incorrectly. The chain of thought here resembles a solution and can be interpreted as such; however, we prefer to term it a chain of thought to emphasize that it replicates a step-by-step reasoning process leading to the answer (additionally, solutions or explanations generally appear \textit{after} the final answer \citep[][\textit{inter alia}]{narang2020wt5,wiegreffe2021reframing,lampinen2022can}).

Chain-of-thought prompting offers several appealing features as a method to enhance reasoning capabilities in language models.

\begin{enumerate}[topsep=1pt,itemsep=0ex]
    \item Firstly, chain of thought fundamentally enables models to break down multi-step problems into intermediate stages, allowing extra computational effort to be devoted to problems that demand more reasoning steps.
    \item Secondly, a chain of thought offers an interpretable insight into the model's behavior, indicating the potential process it used to reach a given answer and providing a means to identify errors in the reasoning path (though fully understanding a model's computations that underpin an answer remains an unresolved challenge).
    \item Thirdly, chain-of-thought reasoning can be applied to tasks such as solving math word problems, performing commonsense reasoning, and carrying out symbolic manipulation, and may, in principle, be extended to any task that humans can solve through language.
    \item Lastly, chain-of-thought reasoning can be easily elicited from sufficiently large pre-trained language models simply by incorporating chain of thought sequence examples into the few-shot prompting exemplars.
\end{enumerate}

Through empirical studies, we will demonstrate the effectiveness of chain-of-thought prompting for arithmetic reasoning (\cref{sec:arithmetic-reasoning}), commonsense reasoning (\cref{sec:commonsense-reasoning}), and symbolic reasoning (\cref{sec:symbolic-reasoning}).

\section{Arithmetic Reasoning}\label{sec:arithmetic-reasoning}
We start by examining math word problems like those illustrated in \cref{fig:pull-figure}, which assess the arithmetic reasoning capabilities of language models. Although straightforward for humans, arithmetic reasoning often poses difficulties for language models \citep[][\textit{inter alia}]{hendrycks2021measuring,patel-etal-2021-nlp}. Remarkably, chain-of-thought prompting, when applied with the 540B parameter language model, achieves performance on par with task-specific finetuned models across several tasks, even setting a new state of the art on the challenging GSM8K benchmark \citep{cobbe2021training}.

\subsection{Experimental Setup}

We investigate the use of chain-of-thought prompting with different language models on several benchmarks.

\textbf{Benchmarks.}
We evaluate on the following five math word problem datasets:
\textbf{(1)} the \textbf{GSM8K} benchmark comprising math word problems \citep{cobbe2021training}, 
\textbf{(2)} the \textbf{SVAMP} dataset containing math word problems with diverse structures \citep{patel-etal-2021-nlp},
\textbf{(3)} the \textbf{ASDiv} dataset featuring a broad range of math word problems \citep{miao-etal-2020-diverse},
\textbf{(4)} the \textbf{AQuA} dataset of algebraic word problems, and
\textbf{(5)} the \textbf{MAWPS} benchmark \citep{koncel-kedziorski-etal-2016-mawps}. Sample problems are provided in Appendix \cref{tab:math-datasets}.

\textbf{Standard prompting.} As a baseline, we use standard few-shot prompting, introduced by \citet{brown2020language}, where a language model is supplied with in-context exemplars consisting of input--output pairs before predicting on a test example. The exemplars are formatted as questions and answers, with the model generating the answer directly, as shown in \cref{fig:pull-figure} (left).

\textbf{Chain-of-thought prompting.} Our method involves enhancing each example in few-shot prompting with a chain of reasoning leading to the corresponding answer, as demonstrated in \cref{fig:pull-figure} (right). Since most datasets only provide an evaluation split, we manually created a set of eight few-shot examples accompanied by chains of thought to use for prompting—\cref{fig:pull-figure} (right) displays one such chain-of-thought example, and the complete set of examples can be found in Appendix \cref{tab:appendix-mwp-prompts}. (These specific examples were not subjected to prompt engineering; robustness analyses are presented in \cref{subsec:robustness} and \cref{subsec:faq-prompt-engineering}.) To determine whether this style of chain-of-thought prompting can reliably trigger effective reasoning across a variety of math word problems, we employed this single set of eight chain-of-thought exemplars across all benchmarks except for AQuA, which involves multiple-choice rather than free-response questions. For AQuA, we utilized four exemplars and their solutions taken from the training set, detailed in Appendix \cref{tab:appendix-aqua-prompt}.

\textbf{Language models.} We conduct evaluations on five large language models. The first is \textbf{GPT-3} \citep{brown2020language}, including the variants text-ada-001, text-babbage-001, text-curie-001, and text-davinci-002, which likely correspond to InstructGPT models with 350M, 1.3B, 6.7B, and 175B parameters respectively \citep{ouyang2022training}. The second model family is \textbf{\lamda{}} \citep{thoppilan2022lamda}, comprising versions with 422M, 2B, 8B, 68B, and 137B parameters. The third is \textbf{\palm{}}, available in sizes of 8B, 62B, and 540B parameters. The fourth model is \textbf{UL2 20B} \citep{tay2022unifying}, and the fifth is \textbf{Codex} \citep[][code-davinci-002 within the OpenAI API]{chen2021evaluating}. For model inference, we use greedy decoding (although subsequent research indicates that chain-of-thought prompting performance can be enhanced by selecting the majority final answer over multiple sampled generations \citep{wang2022self}). For \lamda{}, results are averaged over five random seeds, each corresponding to a distinct random shuffling of the exemplar order. Since \lamda{} experiments showed minimal variation across seeds, we report outcomes for a single exemplar ordering for all other models to reduce computational expense.

\subsection{Results}
The most significant findings of chain-of-thought prompting are summarized in \cref{fig:main-math}, with all experimental outputs across different model families, model sizes, and benchmarks presented in \cref{tab:all-lm-math} in the Appendix. There are three main insights. First, \cref{fig:main-math} demonstrates that chain-of-thought prompting emerges as a capability with increasing model scale \citep{wei2022emergent}. Specifically, chain-of-thought prompting does not enhance performance in smaller models and only begins to provide benefits when applied to models with approximately 100 billion parameters. Our qualitative analysis showed that smaller-scale models generated fluent yet logically flawed chains of thought, resulting in worse performance compared to standard prompting.

\input{main_fables/main-math}
Second, the performance improvements from chain-of-thought prompting are more pronounced on complex tasks. For example, on GSM8K (which has the lowest baseline accuracy), the performance of the largest GPT and \palm{} models more than doubled. Conversely, for SingleOp, the simplest MAWPS subset requiring only a single reasoning step, the improvements were either negligible or negative (see Appendix \cref{tab:mawps-subsets}).

Third, chain-of-thought prompting using GPT-3 175B and \palm{} 540B achieves results comparable to previous state-of-the-art approaches, which generally involve fine-tuning task-specific models on labeled training data. 
As shown in \cref{fig:main-math}, \palm{} 540B with chain-of-thought prompting attains new state-of-the-art performance on GSM8K, SVAMP, and MAWPS (noting that standard prompting already surpassed previous best results on SVAMP). For the remaining two datasets, AQuA and ASDiv, \palm{} employing chain-of-thought prompting comes within 2\% of the state of the art (see Appendix \cref{tab:all-lm-math}).

To gain a deeper understanding of why chain-of-thought prompting is effective, we conducted a manual review of model-generated chains of thought produced by \lamda{} 137B on GSM8K. Among 50 randomly selected instances where the model generated the correct final answer, all but two of the chains of thought were logically and mathematically sound; the exceptions were cases where the chain coincidentally led to the correct answer (refer to \cref{subsec:correct-chain-of-thought} and \cref{tab:appendix-gsm8k-correct-examples} for examples of accurate model-generated chains of thought). We also randomly inspected 50 samples where the model produced incorrect answers. Our analysis revealed that 46\% of these chains were nearly correct, with only minor errors such as calculator mistakes, symbol mapping issues, or a single reasoning step missing, while the remaining 54\% contained significant errors related to semantic understanding or coherence (see \cref{subsec:incorrect-chain-of-thought-analysis}). 

To shed some light on why scaling enhances chain-of-thought reasoning capabilities, we conducted a comparable error analysis for \palm{} 62B and examined whether these errors were rectified by scaling up to \palm{} 540B. The findings indicate that increasing the model size to 540B resolves a substantial portion of errors involving missing single steps and semantic misunderstandings present in the 62B model (see \cref{subsec:why-scale-helps}).

\subsection{Ablation Study}\label{subsec:math-ablation}

The demonstrated advantages of chain-of-thought prompting naturally lead to the question of whether similar performance gains can be achieved through other forms of prompting. \cref{fig:bar_ablation} presents an ablation study exploring three variants of chain-of-thought prompting, described below.

\textbf{Equation only.} One possible reason chain-of-thought prompting is beneficial is that it generates the mathematical equation to be solved. To test this hypothesis, we tried a variant where the model is instructed to output solely the mathematical equation before providing the answer. \cref{fig:bar_ablation} shows that equation-only prompting offers limited improvement for GSM8K, suggesting that the semantic complexity of GSM8K questions makes direct translation to equations challenging without the natural language reasoning steps included in chain-of-thought. However, for datasets consisting of one-step or two-step problems, equation-only prompting does enhance performance, as the equation can be straightforwardly derived from the question (see Appendix \cref{tab:ablations-arithmetic}).

\input{main_fables/ablation-bar}
\textbf{Variable compute only.} Another hypothesis is that chain of thought enables the model to allocate more computation (i.e., intermediate tokens) to more difficult problems. To distinguish the impact of variable computation from chain-of-thought reasoning, we experiment with a setup where the model is prompted to produce a sequence of dots ($\ldots$) matching the number of characters in the equation required to solve the problem. This variant yields performance comparable to the baseline, indicating that variable computation alone does not account for the effectiveness of chain-of-thought prompting, and that expressing intermediate steps in natural language adds value.

\textbf{Chain of thought after answer.} A further possible advantage of chain-of-thought prompting might be that such prompts enable the model to better retrieve relevant knowledge learned during pretraining. To examine this, we test a configuration where the chain of thought prompt is provided only after the answer, isolating whether the model relies on the produced chain of thought to generate the final answer. This variant performs similarly to the baseline, implying that the sequential reasoning captured in the chain of thought serves purposes beyond merely activating prior knowledge.

\input{main_fables/main-robustness}
\subsection{Robustness of Chain of Thought}\label{subsec:robustness}

Sensitivity to exemplar choices is a critical factor in prompting methods—for example, changing the order of few-shot exemplars can cause GPT-3’s accuracy on SST-2 to vary widely, from nearly random chance (54.3\%) to close to state-of-the-art performance (93.4\%) \citep{zhao2021calibrate}. In this concluding subsection, we assess the robustness of chains of thought authored by different annotators. Besides the results above using chains of thought created by Annotator A, two other co-authors of this paper (Annotators B and C) independently composed chains of thought for the same few-shot exemplars (displayed in \cref{sec:appendix-alternate-annotators}). Annotator A also produced an alternative, more concise chain of thought than the original, inspired by the style of solutions in \citet{cobbe2021training}.\footnote{For example, while the original chain of thought uses multiple short sentences (\textit{``There were originally 9 computers. For each of 4 days, 5 more computers were added. So 5 * 4 = 20 computers were added. 9 + 20 is 29.''}), the concise version states \textit{``5 * 4 = 20 new computers were added. So there are 9 + 20 = 29 new computers in the server room now''}.}

\cref{fig:bar-robustness-analysis} presents these findings for \lamda{} 137B applied to GSM8K and MAWPS (additional ablation outcomes for other datasets can be found in Appendix \cref{tab:ablations-arithmetic} / \cref{tab:ablations-commonsense-symbolic}). Although variation exists among different chain of thought annotations, as anticipated when using exemplar-based prompting \citep{le-scao-rush-2021-many,reynolds2021prompt,zhao2021calibrate}, every set of chain of thought prompts significantly outperforms the standard baseline. This indicates that effective chain of thought utilization does not rely on any specific linguistic style.

To verify that effective chain-of-thought prompting generalizes to other exemplar sets, we also conduct experiments using three sets of eight exemplars randomly drawn from the GSM8K training set, which serves as an independent source (examples in this dataset already contain reasoning steps resembling a chain of thought).\footnote{We select examples of length $\leq 60$ tokens to fit within our input context window and restrict examples to $\leq 2$ reasoning steps to ensure a fair comparison with the eight exemplars we created.} \cref{fig:bar-robustness-analysis} demonstrates that these prompts yield performance comparable to our manually crafted exemplars, while also significantly surpassing standard prompting.

Beyond demonstrating robustness to different annotators, independently generated chains of thought, varying exemplars, and diverse language models, we further observe that chain-of-thought prompting for arithmetic reasoning maintains stability across different exemplar orders and varying numbers of exemplars (refer to \cref{subsec:faq-prompt-engineering}).

\section{Commonsense Reasoning}\label{sec:commonsense-reasoning}

While chain of thought is especially well-suited for math word problems, its language-based foundation allows it to be applied broadly to commonsense reasoning tasks, which involve reasoning about physical and human interactions based on assumed general background knowledge. Commonsense reasoning is crucial for effective interaction with the world but remains beyond the capabilities of current natural language understanding systems \citep{talmor2022commonsenseqa}.

\textbf{Benchmarks.}  
We examine five datasets that span a wide array of commonsense reasoning challenges. The well-known \textbf{CSQA} \citep{talmor-etal-2019-commonsenseqa} involves answering commonsense questions about the world that often demand complex semantic understanding and prior knowledge.  
\textbf{StrategyQA} \citep{geva-etal-2021-aristotle} requires models to deduce a multi-step reasoning process to respond to questions. We include two specialized test sets from the BIG-bench initiative \citep{bigbench}: \textbf{Date} Understanding, which entails extracting a date from the provided context, and \textbf{Sports} Understanding, which involves judging whether a sports-related statement is plausible or implausible. Lastly, the \textbf{SayCan} dataset \citep{ahn2022can} focuses on translating natural language instructions into sequences of discrete robot actions.  
\cref{fig:dataset-examples} presents examples with chain of thought annotations across all datasets.

\textbf{Prompts.}  
Our experimental protocol aligns with that of the previous section. For CSQA and StrategyQA, we randomly sampled examples from their training sets and manually crafted chains of thought to serve as few-shot exemplars. Since the two BIG-bench tasks lack training sets, we selected the first ten examples from their evaluation sets as few-shot exemplars and reported results on the remaining examples in the evaluation set. For SayCan, we utilized six training examples from \citet{ahn2022can} and also manually generated chains of thought.

\textbf{Results.}
\cref{fig:commonsense-results} presents these findings for \palm{} (complete results for \lamda{}, GPT-3, and various model sizes can be found in \cref{tab:all-lm-commonsense}).
Across all tasks, increasing model size enhanced the performance of standard prompting; chain-of-thought prompting yielded additional improvements, with the largest gains observed for \palm{} 540B. Using chain-of-thought prompting, \palm{} 540B demonstrated strong results compared to baselines, surpassing the previous state-of-the-art on StrategyQA (75.6\% vs 69.4\%) and exceeding the accuracy of an unaided sports fan on sports understanding (95.4\% vs 84\%).
These findings indicate that chain-of-thought prompting can also boost performance on tasks requiring diverse commonsense reasoning skills (although the improvement was minimal on CSQA).

\input{main_fables/main-commonsense}

\input{main_fables/main-symbolic}
\section{Symbolic Reasoning}\label{sec:symbolic-reasoning}
Our last experimental evaluation focuses on symbolic reasoning, which is straightforward for humans but potentially difficult for language models. We demonstrate that chain-of-thought prompting not only enables language models to solve symbolic reasoning tasks that are hard under standard prompting, but also supports length generalization to inputs at inference time that are longer than the few-shot examples.

\paragraph{Tasks.}
We utilize the following two simple tasks.

\begin{itemize}[leftmargin=*,topsep=0pt]
    \itemsep0em \item \textbf{Last letter concatenation.} This task requires the model to join the last letters of words in a name (e.g., \textit{``Amy Brown''} $\rightarrow$ \textit{``yn''}).
    It is a more difficult variant of first letter concatenation, which language models can already do without chain of thought.\footnote{We tested 10 common names with GPT-3 \texttt{davinci} and it answered all but one correctly.} Full names are generated by randomly combining names from the top one-thousand first and last names according to name census data ({\small{\url{https://namecensus.com/}}}).
    \item \textbf{Coin flip.} This task involves determining whether a coin remains heads up after individuals either flip or do not flip it (e.g., \textit{``A coin is heads up. Phoebe flips the coin. Osvaldo does not flip the coin. Is the coin still heads up?''} $\rightarrow$ \textit{``no''}).
\end{itemize}

Given the well-defined nature of these symbolic reasoning tasks, for each task we evaluate performance on an \textit{in-domain} test set, where the examples involve the same number of steps as the training or few-shot exemplars, as well as on an \textit{out-of-domain} (OOD) test set, where the test examples require more steps than those present in the exemplars. For the last letter concatenation task, the model is exposed only to exemplars consisting of names with two words, and is then tested on concatenation involving names with 3 and 4 words.\footnote{For names longer than two words, multiple first and last names are concatenated.} A similar approach is applied to the number of potential flips in the coin flip task. Our experimental design employs the same methodologies and models as discussed in the preceding sections. We again manually craft chains of thought for the few-shot exemplars in each task, which are displayed in \cref{fig:dataset-examples}.

\paragraph{Results.}
The findings from the in-domain and OOD tests are presented in \cref{fig:symbolic-main} for \palm{}, while results for \lamda{} appear in Appendix \cref{tab:all-lm-symbolic}. Using \palm{} 540B, chain-of-thought prompting achieves nearly perfect solve rates (it is worth noting that standard prompting already succeeds on coin flip with \palm{} 540B, though not with \lamda{} 137B). It is important to recognize that these in-domain evaluations serve as ``toy tasks'' since the ideal solution processes are explicitly provided by the chains of thought in the few-shot exemplars; the model's role is merely to replicate the same steps with new symbols in the test examples. Nevertheless, smaller models fail at this, indicating that the capacity to carry out abstract manipulations on unseen symbols in these three tasks emerges only at model scales around 100 billion parameters.

Regarding the OOD evaluations, standard prompting does not succeed in either task. However, with chain-of-thought prompting, language models display upward scaling trends (although their performance remains lower compared to the in-domain scenario). Therefore, chain-of-thought prompting enables length generalization beyond previously observed chains of thought for language models of adequate size.

\section{Discussion}\label{sec:discussion}

We have investigated chain-of-thought prompting as a straightforward technique to provoke multi-step reasoning behaviors in large language models. Initially, we observed that chain-of-thought prompting significantly enhances performance on arithmetic reasoning tasks, producing gains that surpass those of ablations and remain consistent across various annotators, exemplars, and language models (\cref{sec:arithmetic-reasoning}). Subsequently, experiments focused on commonsense reasoning highlighted how the linguistic character of chain-of-thought reasoning makes it broadly applicable (\cref{sec:commonsense-reasoning}). Lastly, we demonstrated that for symbolic reasoning tasks, chain-of-thought prompting supports OOD generalization to longer sequences (\cref{sec:symbolic-reasoning}). Across all experiments, chain-of-thought reasoning was elicited solely by prompting pre-existing language models, with no fine-tuning conducted during the preparation of this paper.

The appearance of chain-of-thought reasoning as a function of increasing model size has been a recurrent observation \citep{wei2022emergent}. In numerous reasoning tasks where standard prompting shows a flat scaling curve, chain-of-thought prompting results in sharply rising scaling curves. Chain-of-thought prompting seems to broaden the range of tasks that large language models can successfully address—in other words, our findings emphasize that standard prompting only establishes a lower bound on the capabilities of large language models. This insight likely provokes more questions than it resolves—for example, how much further might reasoning skills improve with additional model scaling? What alternative prompting techniques could expand the set of tasks solvable by language models?

Regarding limitations, we first clarify that although chain-of-thought mimics the reasoning processes of humans, this does not confirm whether the neural network is genuinely “reasoning,” which we leave as an open question. Second, while the manual addition of chains of thought to exemplars incurs minimal cost in a few-shot setting, such annotation expenses could be prohibitive for fine-tuning (although this might be overcome through synthetic data generation or zero-shot generalization). \orange{Third, there is no assurance that reasoning paths are correct, which can result in both accurate and inaccurate answers; enhancing the factual accuracy of language model generations remains a direction for future research \citep[][\textit{inter alia}]{rashkin2021measuring,ye2022unreliability,wiegreffe2021reframing}.} Lastly, the fact that chain-of-thought reasoning emerges only at large model scales makes deployment costly in practical applications; future work could investigate methods to induce reasoning capabilities in smaller models.

\section{Related Work} This research draws inspiration from multiple fields, which we discuss in more detail in an extended related work section (\cref{sec:extended-related-work}). Here, we focus on two key directions and their associated studies that are arguably the most pertinent.

The first relevant direction involves employing intermediate steps to tackle reasoning problems. \cite{ling-etal-2017-program} introduce the concept of utilizing natural language rationales to address math word problems by breaking them down into a sequence of intermediate steps. Their approach stands in notable contrast to prior work that uses formal languages for reasoning \citep{roy-2016-reasoning, chiang-chen-2019-semantically, amini-etal-2019-mathqa, chen2019neural}. Building on \citet{ling-etal-2017-program}, \citet{cobbe2021training} create a larger dataset and use it to finetune a pretrained language model rather than training one from scratch. In program synthesis, \cite{nye2021show} utilize language models to forecast the final outputs of Python programs by first predicting intermediate computational steps line-by-line, demonstrating that this stepwise prediction technique surpasses directly predicting the final results.

This paper is also closely connected to the extensive recent literature on prompting. Since the introduction of few-shot prompting by \citet{brown2020language}, various general strategies have been proposed to enhance models' prompting capabilities, such as automatically learning prompts \citep{lester-etal-2021-power} or providing models with task-descriptive instructions \citep{wei2021finetuned,sanh2021multitask,ouyang2022training}. While these methods improve or augment the input portion of the prompt (e.g., instructions prepended to inputs), our approach takes an orthogonal path by enhancing the outputs of language models through the generation of a chain of thought.

\section{Conclusions}
We have investigated chain-of-thought prompting as a straightforward and widely applicable technique to improve reasoning in language models. Across experiments involving arithmetic, symbolic, and commonsense reasoning, we observe that chain-of-thought reasoning emerges as a function of model scale, enabling sufficiently large language models to succeed on reasoning tasks that typically exhibit flat scaling curves. Expanding the scope of reasoning tasks solvable by language models may encourage further research on language-based reasoning methods.

\end{document}